\section{Arquitetura dos Experimentos}
\label{sec:arquitetura-experimentos}

Para avaliar o sistema em condições distintas, foi criada uma camada experimental em
torno do \textit{SchedulingEngine}. Esta camada separa a definição dos cenários da sua
execução, o que permite variar sistematicamente a dimensão do problema, o número de
processos, o limite de qubits, os parâmetros do QAOA e a estratégia de resolução sem
alterar o código principal do sistema.

A arquitetura experimental organiza-se em quatro etapas: configuração do cenário,
orquestração da execução, resolução pela pipeline selecionada e persistência dos
resultados. A Figura~\ref{fig:fluxo-experimentos} apresenta este fluxo. Os executores de
cenários constituem a camada adicional sobre a arquitetura descrita no capítulo anterior;
o restante processamento continua a ser realizado pelo \textit{SchedulingEngine}.

\begin{figure}[H]
\centering
\begin{tikzpicture}[
    node distance=1.4cm,
    box/.style={rectangle, draw, rounded corners, align=center, minimum width=3.7cm, minimum height=0.9cm},
    arrow/.style={->, thick}
]

\node[box] (toml) {Ficheiro TOML\\do cenário};
\node[box, below of=toml] (runner) {Scenario Runner\\ou Sweep Runner};
\node[box, below of=runner] (engine) {Scheduling Engine};
\node[box, below left=1.2cm and 0.8cm of engine] (default) {Pipeline direta};
\node[box, below right=1.2cm and 0.8cm of engine] (iterative) {Pipeline iterativa\\com sub-QUBOs};
\node[box, below=2.4cm of engine] (solver) {Solver QAOA\\e validação};
\node[box, below of=solver] (results) {Resultados\\JSON/JSONL, TXT e logs};

\draw[arrow] (toml) -- (runner);
\draw[arrow] (runner) -- (engine);
\draw[arrow] (engine) -- (default);
\draw[arrow] (engine) -- (iterative);
\draw[arrow] (default) -- (solver);
\draw[arrow] (iterative) -- (solver);
\draw[arrow] (solver) -- (results);

\end{tikzpicture}
\caption{Arquitetura utilizada na execução dos experimentos da investigação.}
\label{fig:fluxo-experimentos}
\end{figure}

\subsection{Definição e Execução dos Cenários}

Cada experimento é descrito num ficheiro de configuração em formato TOML. As secções
do ficheiro definem o \textit{workload}, a formulação QUBO, os parâmetros do QAOA, a
decomposição, a validação e as opções de execução. Esta separação permite repetir uma
experiência ou alterar uma única variável experimental sem modificar o código do
sistema.

O \textit{Scenario Runner} carrega a configuração, converte os seus valores nos objetos
usados pelo sistema e executa o número de repetições especificado. Nos estudos
paramétricos, o \textit{Sweep Runner} expande os eixos e casos definidos no ficheiro TOML,
gera as combinações correspondentes e entrega cada variante ao mesmo mecanismo de
execução. Deste modo, tanto uma execução isolada como uma varredura de parâmetros seguem
o mesmo percurso e produzem registos com a mesma estrutura.

Para cada variante, o \textit{SchedulingEngine} seleciona a pipeline direta ou a pipeline
iterativa. A primeira resolve a instância QUBO completa; a segunda decompõe o problema em
sub-QUBOs quando o limite configurado de qubits é ultrapassado. A solução produzida pelo
solver é depois validada e resumida num conjunto comum de métricas.

Um exemplo simplificado de um cenário experimental é apresentado de seguida:

\begin{lstlisting}[caption={Exemplo simplificado de um cenário experimental em TOML.}]
scenario_version = 1
id = "T3.5"
name = "Three-core iterative decomposition"

[workload]
mode = "preset"
num_processes = 5
num_cores = 3
weights = [0.36, 0.27, 0.22, 0.18, 0.12]

[qubo]
penalty = 5.0

[qaoa]
layers = 1
steps = 50
learning_rate = 0.05
top_k = 20
mixer_type = "xy"

[decomposition]
qubit_max = 9
num_cores = 3
sorting_strategy = "WEIGHT_DESCENDING"

[execution]
pipeline = "iterative"
solvers = ["qaoa"]
repeats = 1
\end{lstlisting}

\subsection{Recolha e Armazenamento dos Resultados}

Após cada execução, o sistema cria uma pasta própria para o cenário e guarda uma
cópia da configuração utilizada. Cada repetição origina um ficheiro JSON com a
configuração resolvida e as métricas, um ficheiro TXT legível e um log da saída produzida.

Em cada execução são registados, entre outros, os seguintes dados:

\begin{itemize}
    \item configuração usada no experimento;
    \item número de processos, núcleos e variáveis QUBO;
    \item energia da solução obtida;
    \item viabilidade da solução;
    \item tempo de resolução;
    \item atribuição final de processos a cores;
    \item resultado da validação;
    \item desequilíbrio final de carga entre os núcleos;
\end{itemize}
